\scanpage{16}
\mspage{13}

Il est clair que $\unsure{\mathcal A|\mathscr A}_{I}\overset{\mathrm{d\acute ef}}{=}\mathbb D_{\unsure{\mathcal M|\mathscr M},\Sigma}(I)$ dépend de $I$ de façon fonctorielle contravariante, et $=$ de façon 2-fonctorielle.

Bien entendu, pour avoir vraiment un dérivateur, il faut vérifier les axiomes. Le plus fort est l'existence des $f_{!},f_{*}$. Elle est prouvée par \uline{\unsure{Thomason|Thomasson}} lorsque $(\mathcal M,\Sigma)$ provient d'une ``catégorie de \unsure{modèles|modèle}'' au sens de Quillen. Ce n'est pas clair pour moi qu'alors $\mathcal M_{I},\Sigma_{I}$ provienne également d'une catégorie de modèles (dont p. ex. les fibrations et les cofibrations seraient définies, à partir de ces données dans $\mathcal M$, comme tantôt pour les équivalences).

Quand on prend pour $(\mathcal M,\Sigma)$ une catégorie de modèles pour $\mathrm{Hot}$ (p. ex les ens. $1/2$ \unsure{simpliciaux}, \[
\begin{aligned}
\mathcal M&=\underline{\mathrm{Hom}}\,(\unsure{\mathbb{\Delta}|\Delta}^{\unsure{\circ|o}},\mathrm{Ens}),\quad\text{d'où}\quad\mathcal M_{I}\simeq\underline{\mathrm{Hom}}\,(\mathbb{\Delta}^{\circ}\times I^{\circ},\mathrm{Ens})\\
&\simeq\underline{\mathrm{Hom}}\,(\mathbb{\Delta}^{\circ},\underbrace{\underline{\mathrm{Hom}}\,(I^{\circ},\mathrm{Ens})}_{I^{\wedge}}),
\end{aligned}
\] \inserted{cat des} objets \unsure{semisimpliciaux|semi-simpliciaux} \unsure{dans|du} le topos $I^{\wedge}$,), on trouve une description d'un dérivateur (non additif), qui jouera un rôle particulièrement important. Je le désigne par

\[
\mathrm{HOT},
\]

sa catégorie de coeff. fondamentale est $\mathrm{Hot}$, la catégorie homotopique standart.

\figure{\begin{tikzpicture}\draw (0,.34) -- (.25,0);\end{tikzpicture}}

% archive sheet number: 15
